\begin{description}
	\item[Estimation for knowing the population parameters] In statistics and psychological statistics in university, inferential statistics are used to find numbers that represent the population, instead of just examining descriptive statistics of samples.
	\item[Inference using probability] There are three ways to know the population parameters: the moment method using representative values, maximum likelihood method using probability distributions, and the Bayesian method.
	\item[Experimental design and linear models] In psychology, research conditions are controlled using experimental design, and discussion is often based on the difference in the mean values (average causal effect). These can generally be expressed as linear models.
	\item[Model comparison and decision making] There are methods such as model comparison and NHST to not only estimate population parameters but also reach certain "conclusions" or "decisions" based on them.
	\item[Practical use of statistical programming language R] Learning to use the R programming language to not only understand the theory but also actually perform calculations is a necessary skill.
\end{description}
